% ORGANIZACIÓN DE COMPUTADORAS
% SEGUNDO CUATRIMESTRE 2017
% SEGUNDO EXAMEN PARCIAL

\documentclass[12pt,a4paper]{article}
\usepackage[spanish]{babel}
\selectlanguage{spanish}
\usepackage[utf8]{inputenc}

\usepackage{graphicx}
\usepackage{multirow}
\usepackage{xcolor}
\usepackage{fancyhdr}
\renewcommand{\headrulewidth}{0pt}

\setlength{\textwidth}{170mm}
\setlength{\oddsidemargin}{-5.4mm}
\setlength{\evensidemargin}{-5.4mm}
\setlength{\textheight}{247mm}
\setlength{\topmargin}{-15.4mm}
\setlength{\headheight}{11mm}
\setlength{\headsep}{2mm}

% CARPETA CON IMAGENES DE ENCABEZADO
\graphicspath{{estilo/}}

% DEFINICIÓN DE COMANDOS

\newcommand{\setyear}[1]{\year=#1}
\newcommand{\universidad}{Universidad Nacional del Sur}
\newcommand{\departamento}{Departamento de Ciencias e Ingeniería de la Computación}
\newcommand{\materia}{Organización de computadoras}
\newcommand{\cuatrimestre}{Segundo}

\newcommand{\RISC}{\textsc{risc}}
\newcommand{\CISC}{\textsc{cisc}}
\newcommand{\OCUNS}{\textsc{ocuns}}
\newcommand{\PC}{\textsf{PC}}
\newcommand{\SP}{\textsf{SP}}

% ENCABEZADO GENERAL
% DeptoLogo - Departamento - Universidad - Cuatrimestre - Año - UniLogo
\newcommand{\header}{
	\begin{center}
		\begin{tabular}{lcl}
		\multirow{4}{*}
		{\smash{$\vcenter{\hbox{\includegraphics[width=.78in]{DeptoLogo}}}$}} & {\Large\textsc{\materia}} &  
		\multirow{4}{*}
		{\smash{$\vcenter{\hbox{\includegraphics[width=.88in]{UniLogo}}}$}} \\
		         & \departamento & \\
		         & \universidad & \\
		         & \textbf{\cuatrimestre\ Cuatrimestre de \number\year} & \\
		\end{tabular}\\[2em]
	\end{center}
}

% ENCABEZADO EDITABLE%
% Uso: \Titulo{Primer texto negrita}{Segundo texto normal}
% header + Texto 1 + Texto 2
\newcommand{\Titulo}[2]{ 
	\begin{center}
		\header
		 \textbf{#1} \\ #2 \\[4pt]
	\end{center}
}

% ENCABEZADO PRÁCTICOS CON NUMEROS
% Uso: \TPnumero{Número de Práctico}{Título de Práctico}
% header + Trabajo Práctico Nº X + Titulo de Práctico
\newcommand{\TPnumero}[2]{ 
	\begin{center}
		\header
        Trabajo Práctico N$^{\circ}$ #1 \\ \textbf{#2} \\[4pt]
    \end{center}
}

% ENCABEZADO PRÁCTICOS CON LETRAS
% Uso: \TPletra{Letra de Práctico}{Título de Práctico}
% header + Trabajo Práctico 'X' + Titulo de Práctico
\newcommand{\TPletra}[2]{ 
	\begin{center}
		\header
		Trabajo Práctico #1 \\ \textbf{#2} \\[4pt]
	\end{center}
}

% ENCABEZADO PRÁCTICOS PROGRAMACIÓN CON NUMEROS
% Uso: \TPPnumero{Número de Práctico}{Título de Práctico}
% header + Trabajo Práctico de Programación Nº X + Titulo de Práctico
\newcommand{\TPPnumero}[2]{ 
	\begin{center}
		\header
		Trabajo Práctico de Programación N$^{\circ}$ #1 \\ \textbf{#2} \\[4pt]
	\end{center}
}

% ENCABEZADO PRÁCTICOS PROGRAMACIÓN CON LETRAS
% Uso: \TPPletra{Letra de Práctico}{Título de Práctico}
% header + Trabajo Práctico de Programación 'X' + Titulo de Práctico
\newcommand{\TPPletra}[2]{ 
	\begin{center}
		\header
		Trabajo Práctico de Programación #1 \\ \textbf{#2} \\[4pt]
	\end{center}
}

% ENCABEZADO PROYECTO CON NÚMERO
% Uso: \PRnumero{Número de Proyecto}{Título de Proyecto}
% header + Proyecto Nº X + Titulo de Proyecto
\newcommand{\PRnumero}[2]{
	\begin{center}
		\header
		Proyecto N$^{\circ}$ #1 \\ \textbf{#2} \\[4pt]
	\end{center}
}

\begin{document}

\Examen{Segundo Examen Parcial}

\begin{centering}
\emph{Apague cualquier dispositivo electrónico en su poder y manténgalo guardado. No puede utilizar auriculares. Lea todo el ejercicio antes de comenzar a desarrollarlo.}	
\end{centering}

\PEjercicio{1} Hacer

\PEjercicio{2} Hacer

\PEjercicio{3} Hacer

\PEjercicio{4} Hacer

\PEjercicio{5} Considerando el siguiente programa para la arquitectura OCUNS, en la que toda lectura/escritura sobre la dirección FF es redireccionada a la E/S estándar: \\ [2.5mm]
\begin{small}
	\begin{minipage}{.4\textwidth}
		\begin{verbatim}
			      LDA R0, FFh
			      LOAD R1, 0(R0)
			      LOAD R2, 0(R0)
			      XOR R3, R3, R3
			      LDA R4, lbl3
			      JZ R1, lbl3
			      JZ R2, lbl3
			      SUB R5, R1, R2
			      JG R5, lbl2
			lbl1: ADD R3, R3, R2
			      DEC R1
			      JZ R1, lbl3
			      JMP R4
			lbl2: ADD R3, R3, R1
			      DEC R2
			      JZ R2, lbl2
			lbl3: STORE R3, 0(R0)
			      HLT
		\end{verbatim}
	\end{minipage}
	\begin{minipage}{.6\textwidth}	
		\begin{tabular}{|c|l|c|l|}\hline
			\textsc{Op.} & \textsc{Descr.} & \textsc{Form.} & \textsc{Pseudocódigo} \\ \hline
			\textbf{0} & \textsf{add} & \textbf{I} & \texttt{R[d] $\leftarrow$ R[s] + R[t]} \\
			\textbf{1} & \textsf{sub} & \textbf{I} & \texttt{R[d] $\leftarrow$ R[s] - R[t]} \\
			\textbf{2} & \textsf{and} & \textbf{I} & \texttt{R[d] $\leftarrow$ R[s] \& R[t]} \\
			\textbf{3} & \textsf{xor} & \textbf{I} & \texttt{R[d] $\leftarrow$ R[s] \^{} R[t]} \\
			\textbf{4} & \textsf{lsh} & \textbf{I} & \texttt{R[d] $\leftarrow$ R[s] <{}<{} R[t]} \\
			\textbf{5} & \textsf{rsh}  & \textbf{I} & \texttt{R[d] $\leftarrow$ R[s] >{}>{} R[t]} \\
			\textbf{6} & \textsf{load}  & \textbf{I} & \texttt{R[d] $\leftarrow$ mem[offset + R[s]]} \\
			\textbf{7} & \textsf{store} & \textbf{I} & \texttt{mem[offset + R[d]] $\leftarrow$ R[s]} \\
			\textbf{8} & \textsf{lda}   & \textbf{II} & \texttt{R[d] $\leftarrow$ addr} \\
			\textbf{9} & \textsf{jz}    & \textbf{II} & \texttt{if (R[d] == 0) PC $\leftarrow$ PC + addr} \\
			\textbf{A} & \textsf{jg}    & \textbf{II} & \texttt{if (R[d] >{} 0) PC $\leftarrow$ PC + addr} \\
			\textbf{B} & \textsf{call}  & \textbf{II} & \texttt{R[d] $\leftarrow$ PC; PC $\leftarrow$ addr} \\
			\textbf{C} & \textsf{jmp}   & \textbf{III} & \texttt{PC $\leftarrow$ R[d]} \\
			\textbf{D} & \textsf{inc}   & \textbf{III} & \texttt{R[d] $\leftarrow$ R[d] + 1} \\
			\textbf{E} & \textsf{dec}   & \textbf{III} & \texttt{R[d] $\leftarrow$ R[d] - 1} \\
			\textbf{F} & \textsf{hlt}   & \textbf{III} & \texttt{exit} \\ \hline
		\end{tabular}	
	\end{minipage}
	\begin{center}
		\begin{tabular}{*{17}{c}}
			\textsf{FORMATO} & 15 & 14 & 13 & 12 & 11 & 10 & 9 & 8 & 7 & 6 & 5 & 4 & 3 & 2 & 1 & 0 \\ \hline
			\multicolumn{1}{|c|}{\textbf{I}} & 0 & $\times$ & $\times$ & $\times$ &
			\multicolumn{4}{|c|}{\textsf{dest. d}} &
			\multicolumn{4}{|c|}{\textsf{src. s}} &
			\multicolumn{4}{|c|}{\textsf{src. t / off.}} \\ \hline
			\multicolumn{1}{|c|}{\textbf{II}} & 1 & 0 & $\times$ & $\times$ &
			\multicolumn{4}{|c|}{\textsf{dest. d}} &
			\multicolumn{8}{|c|}{\textsf{address addr}} \\ \hline
			\multicolumn{1}{|c|}{\textbf{III}} & 1 & 1 & $\times$ & $\times$ &
			\multicolumn{4}{|c|}{\textsf{dest. d}} &
			\multicolumn{8}{|c|}{\textsf{-}} \\ \hline
		\end{tabular}
	\end{center}		
\end{small}

\begin{enumerate}[a)]
	\item Ensamblar el programa a partir de la dirección 00h. \label{ensamblado}
	\item Si se reubicara el código máquina obtenido en el inciso (\ref{ensamblado}) a partir de la dirección 20h, ¿qué referencias a memoria requieren ser ajustadas? Justificar adecuadamente.
	\item Suponiendo que los valores ingresados por teclado son 1Ah y 04h:
	\begin{enumerate}[c.1)]
		\item Realice una traza mostrando la evolución del contenido de cada registro.
		\item Describir el propósito del programa en su conjunto.
		\item ¿Qué sucede con el resultado retornado por el programa si los valores ingresados fueran 04h y 1Ah? ¿Cuál es la diferencia?
		\item ¿Cuál es la restricción en cuanto a los valores de entrada que considera el programa para su correcto funcionamiento? 
	\end{enumerate}
\end{enumerate}

\end{document}


